\section{Discussion}
\label{sec:discussion}
In this paper, we focus on the goodput-optimized setting and propose \sys under the large-scale LLM serving scenario. As LLMs are widely used and deployed across various service scenarios with different optimization targets and resource limits, it becomes almost impossible to find a one-size-fits-all solution that effectively addresses all aspects of LLM serving. In this section, we discuss the pros and cons of \sys and potentially better solutions in other scenarios.

\parabf{Throughput-optimized scenarios.} In offline applications that are not latency-sensitive, users typically have lower requirements for response time~\cite{sheng2023flexgen}. This allows serving systems to shift focus towards maximizing overall throughput instead of goodput and the effectiveness of DistServe may be compromised. In this case, techniques such as chunked-prefill with piggyback~\cite{agrawal2023sarathi,deepspeed_mii} may be preferred since it can fill each batch to the compute-bound threshold, thereby maintaining higher GPU utilization in every iteration.

\parabf{Resource-constrained scenarios.} Small-scale enterprises and individual researchers often lack the resources to deploy LLMs on large-scale clusters~\cite{sheng2023flexgen, song2023powerinfer}. In resource-constrained scenarios, such as environments with only a few or even a single GPU, the design space for DistServe is significantly limited. It struggles or even fails to adjust the parallel strategies and resource allocation to effectively enhance serving performance. In this case, simpler architectural choices like non-disaggregated systems~\cite{vllm, deepspeed_mii} may reduce deployment complexity and optimize operational efficiency.

\parabf{Long-context scenarios.} Nowadays, more and more GPT models support extremely long contexts, such as Claude-3~\cite{claude3}, Gemini-1.5~\cite{gemini}, and Large World Model (LWM)~\cite{liu2023world}, which all have a 1M context window. In such scenarios, the transmission overhead will increase as the size of the KV cache grows linearly with the prompt length. However, the prefill computation grows quadratically, so the relative duration of transmission and prefill job decreases. Meanwhile, a longer context further exacerbates the disparity in computational demands between prefill and decoding jobs, leading to increased interference between them. Therefore, the disaggregation approach proposed in \sys remains promising in long-context serving.